\chapter{Contesto aziendale e tecnologico}
\label{chap:contesto-aziendale}

Questo capitolo presenta l'organizzazione ospitante dal punto di vista delle tecnologie utilizzate, dei processi interni di sviluppo e manutenzione, della tipologia di clientela e della propensione all'innovazione, sulla base di osservazione diretta durante lo stage.

\section{L'azienda}
\label{sec:azienda}

% TODO: descrizione dell'organizzazione ospitante (settore, dimensione,
% posizionamento sul mercato), con indicazione esplicita di come l'informazione
% riportata è stata acquisita (colloqui, osservazione diretta, documentazione
% interna) e di quanto essa sia da ritenersi accurata. Evitare toni promozionali.


\section{Tecnologie e processi di sviluppo}
\label{sec:tecnologie-processi}

% TODO: stack tecnologico adottato dall'azienda; processi interni di sviluppo,
% manutenzione e organizzazione del lavoro (metodologie, strumenti di gestione
% del codice, ciclo di rilascio) osservati durante lo stage.


\section{Clientela e propensione all'innovazione}
\label{sec:clientela-innovazione}

% TODO: tipo di clientela servita (dimensione, settore pubblico/privato) e
% rapporto storico dell'azienda con l'innovazione tecnologica, come premessa
% al capitolo successivo sul ruolo degli stage in questo processo.
